\newpage
\section{Conclusioni}
In sintesi, il tirocinio al Fablab è stata un'esperienza estremamente pratica che mi ha permesso di trasformare la teoria in progetti reali. Dal punto di vista tecnico, ho acquisito un'autonomia completa: dalla modellazione CAD alla fabbricazione digitale (stampa 3D e laser), integrando l'elettronica e la programmazione per dare vita a prototipi complessi. Anche lo sviluppo di script personali per automatizzare il mio lavoro mi ha dimostrato quanto l'informatica possa semplificare e velocizzare le attività quotidiane.

Altrettanto importante è stata la mia crescita personale. Gestire i workshop per studenti internazionali ha migliorato notevolmente il mio inglese tecnico e la mia capacità di spiegare concetti complessi in modo semplice. Ho imparato a non bloccarmi davanti a un imprevisto o a una macchina guasta, ma a cercare la soluzione pratica per arrivare al risultato. Esco da questo percorso con una forte sicurezza nelle mie capacità, portando con me un bagaglio di competenze tecniche e relazionali che sarà fondamentale per il mio futuro.